% Authoritative module; see specific_highlights_index.md.
\ifdefined\HighlightsRendererLoaded\endinput\fi
\let\HighlightsRendererLoaded\relax
\RequirePackage{tikz}
\newif\ifGTenHighlights
\GTenHighlightstrue
% Switch off with \GTenHighlightsfalse (restore with \GTenHighlightstrue).
% Use \HighlightStroke{color}{fragment} with any xcolor color.
% Every role command accepts ONE short text OR math fragment, in its current
% mode: \GTenHLValue{buyers}, $\GTenHLValue{X}_i^2$,
% \[p_X(\GTenHLValue{1})=\GTenHLProbability{0.5}.\]
% In aligned equations keep &, \\, relation signs and binary operators OUTSIDE
% strokes; highlight atoms/short phrases, not whole calculations. Do not nest
% strokes. For text within math use \text, e.g. \GTenHLOutcome{\text{cash}}.
% mathpalette preserves display/text/script/scriptscript sizes without nesting
% an equation environment. The node's base is exactly the content's baseline.
% Only the background is translucent; glyphs retain their current color.
\tikzset{gten marker/.style={rectangle,rounded corners=1mm,draw=none,fill=#1,
  fill opacity=.32,text opacity=1,inner xsep=1.5mm,inner ysep=0.8mm,
  outer sep=0pt}}
% Private implementation avoids LaTeX @ names so this input needs no catcode
% changes. \mathsurround is reset locally inside each math node.
\newcommand{\GTenHLMathStroke}[2]{%
  \tikz[baseline=(gtenhl.base)]\node[gten marker=\GTenHLCurrentColor]
    (gtenhl) {$\mathsurround=0pt #1#2$};}
\newcommand{\HighlightStroke}[2]{%
  \ifGTenHighlights
    \begingroup\def\GTenHLCurrentColor{#1}%
    \ifmmode\mathpalette\GTenHLMathStroke{#2}%
    \else\tikz[baseline=(gtenhl.base)]\node[gten marker=#1]
      (gtenhl) {#2};\fi
    \endgroup
  \else#2\fi}
\let\GTenHLStroke\HighlightStroke

